\documentclass{article}
\usepackage[utf8]{inputenc}
\usepackage{amsmath}
\usepackage{amssymb}

\begin{document}


% -> Portada
\begin{titlepage}
    \centering
    \vspace*{5cm} 

    {\Huge \bfseries Entrega 1 \par}
    
    \vspace{1.5cm}

    {\Large Cristóbal Cabrera \par}
    {\Large Verónica Torres \par}    
    {\Large Alfonso Pavez \par}    
    \vspace{2cm} 
    
    % Información adicional
    {\large Optimización \par}
    {\large Yoslandi Lazo \par}
    {\large Universidad Técnica Federico Santa María\par}
    
    \vspace{1.5cm}
    
    {\large \today \par} % \today inserta la fecha actual
    \vfill 
    
\end{titlepage}


% -> Modelo matematico
\section*{Formulación del Modelo Matemático}

\subsection*{Conjuntos e Índices}
\begin{itemize}
    \item $I$: Conjunto de trabajadores, indexado por $i$.
    \item $H$: Conjunto de días del horizonte de planificación, indexado por $h$.
    \item $T$: Conjunto de turnos diarios (ej. $T = \{\text{mañana, tarde, noche}\}$), indexado por $t$.
    \item $F$: Subconjunto de los días $H$ que corresponden a fines de semana, indexado por $f$.
\end{itemize}

\subsection*{Parámetros}
\begin{itemize}
    \item $D_{iht}$: Disposición del trabajador $i$ para realizar el turno $t$ en el día $h$, donde $D_{iht} \in \{0, 1, ..., 10\}$.
    \item $R_{ht}$: Número requerido de trabajadores para el turno $t$ en el día $h$, donde $R_{ht} \in \mathbb{N}$.
\end{itemize}

\subsection*{Variables de Decisión}
\begin{itemize}
    \item $X_{iht} \in \{0, 1\}$: Variable binaria que indica si al trabajador $i$ se le asigna el turno $t$ en el día $h$.
    \item $W_{if} \in \{0, 1\}$: Variable binaria auxiliar que indica si el trabajador $i$ trabaja en el fin de semana $f$.
\end{itemize}

\subsection*{Función Objetivo}
El objetivo es maximizar la disposición total del personal:
\begin{equation}
    \text{Maximizar} \quad Z = \sum_{i \in I} \sum_{h \in H} \sum_{t \in T} D_{iht} \cdot X_{iht}
\end{equation}

\subsection*{Restricciones}
\begin{enumerate}
    \item \textbf{Naturaleza de las Variables:} Define el dominio de las variables.
    \begin{align}
        X_{iht} &\in \{0, 1\} \quad \forall i \in I, \forall h \in H, \forall t \in T \\
        W_{if} &\in \{0, 1\} \quad \forall i \in I, \forall f \in F
    \end{align}

    \item \textbf{Cobertura de demanda: }
    Cada turno en cada dia debe ser cubierto por almenos el numero requerido de trabajadores.
    \begin{equation}
        \sum_{i \in I} X_{iht} \ge R_{ht}
        \quad \forall h \in H, \forall t \in T
    \end{equation}

    \item \textbf{Disponibilidad: }
    Un trabajador solo puede ser asigando turnos para los cuales su disponibilidad sea mayor a 0.

    \begin{equation}
        X_{iht} \le
        \begin{cases}
        1, & \text{si } D_{iht} > 0\\
        0, & \text{si } D_{iht} = 0
        \end{cases}
        \quad \forall i \in I, \forall h \in H, \forall t \in T
    \end{equation}

    \item \textbf{Maximo de dos turnos por dia:}
    Un trabajador no puede realizar dos turnos el mismo dia
    \begin{equation}
        \sum_{t \in T} X_{iht} \le 2
        \quad \forall i \in I, \forall h \in H
    \end{equation}
    \item \textbf{Prohibicion de turno noche seguido de turno mañana:}
    Si un trabajador realiza el turno de noche en el dia $h$, este no podra realizar el turno de mañana. 
    \begin{equation}
        X_{ih, \text{noche}} + X_{ih, \text{mañana}} \le 1
        \quad \forall i \in I, \forall h \in H \setminus \{\text{ultimo dia}\}
    \end{equation}
    \item \textbf{Definicion de fin de semana trabajado:}
    $W_{if}$ se activa si el trabajador $i$ realiza algún turno durante el fin de semana $f$.
    \begin{align}
        W_{if} &\ge X_{iht} && \forall i \in I, \forall f \in F, \forall h \in f, \forall t \in T \\
        W_{if} &\le \sum_{h \in f} \sum_{t \in T} X_{iht} && \forall i \in I, \forall f \in F
    \end{align}
    \item \textbf{No puede trabajar tres fin de semanas consecutivos}
    Ningun trabajador puede trabajar tres fin de semanas seguidos
    \begin{equation}
        W_{i,f} + W_{i, f+1} + W_{i, f+2} \le 2
        \quad \forall i \in I, \forall f = 1, \dots |F| -2
    \end{equation}
    
    \item \textbf{Dominio de las variables:} 
    \begin{align}
        X_{iht} &\in \{0,1\} && \forall i \in I, \forall h \in H, \forall t \in T\\
        W_{if} &\in \{0,1\} && \forall i \in I, \forall f \in F
    \end{align}
    
\end{enumerate}

\end{document}